% Prospectos Solares. Presentacion comercial.
%
% Version comercial. Vende el SERVICIO de prospeccion, no los lotes de una
% corrida concreta. Publico: un desarrollador o un inversionista que conoce el
% sector. Que sepa que es la irradiacion se da por hecho.
%
% Comparte estilo.tex y las figuras con la version tecnica, de modo que las dos
% se ven como un mismo producto y una correccion de sistema alcanza a las dos.
%
% Regla de redaccion de ESTA version: nada de umbrales, percentiles, nombres de
% capa, recuentos de una corrida ni jerga de metodo. Y poco texto: la lamina
% ensena, el expositor cuenta. Lo que el cliente necesita saber es que problema
% se le resuelve, como se le resuelve y que recibe.
%
% Compilar desde presentacion_sb/:
%   pdflatex -output-directory=../outputs/presentacion_sb/build prospectos_solares_comercial.tex
%   pdflatex -output-directory=../outputs/presentacion_sb/build prospectos_solares_comercial.tex
%   python verificar.py ../outputs/presentacion_sb/build/prospectos_solares_comercial.pdf

\documentclass[aspectratio=169, 11pt]{beamer}

\usepackage[utf8]{inputenc}
\usepackage[T1]{fontenc}
\usepackage[spanish, es-nodecimaldot]{babel}
\usepackage{graphicx}
\usepackage{booktabs}
\usepackage{array}
\usepackage{tikz}
\usetikzlibrary{arrows.meta, calc}
\usepackage{xcolor}
\usepackage{helvet}
\renewcommand{\familydefault}{\sfdefault}
\renewcommand{\arraystretch}{1.25}

\graphicspath{%
  {../outputs/presentacion_sb/figuras/}%
  {../outputs/presentacion_sb/capturas/}}
% ============================================================================
% Prospectos Solares. Sistema visual de la presentacion.
% ----------------------------------------------------------------------------
% La paleta, el isotipo y las proporciones se toman del reporte de lotes
% (reporte_predios/plantilla.py, modo claro), para que la presentacion y el
% entregable se lean como un mismo producto y no como dos piezas distintas.
%
% Reglas del sistema:
%   Color      un solo verde azulado de marca. El acento oscuro (acento) es el
%              que lleva texto y linea; el claro (marca) solo identifica.
%              bien, aviso y mal se reservan para estado, nunca para decorar.
%   Rejilla    margen lateral unico (\margenlat) y altura util unica
%              (\altoutil). Ninguna lamina cambia de proporcion.
%   Filete     un solo divisor de marca: 1,8 cm por 1,4 pt en acento. Las
%              separaciones internas usan la hairline de 0,6 pt en linea.
%   Tipografia una sola familia de palo seco, tres pesos de gris (tinta,
%              tinta2, tinta3) y jerarquia por tamano, no por color.
%   Cajas      toda caja de diagrama es un nodo de tikz con text width y
%              minimum height: la caja se dimensiona a partir del texto, de
%              modo que el texto nunca puede desbordarla.
%   Lectura    sin justificar y sin partir palabras.
% ============================================================================

\usetheme{default}
\usefonttheme{professionalfonts}
\setbeamertemplate{navigation symbols}{}

% --- Paleta ------------------------------------------------------------------
\definecolor{marca}{HTML}{18919C}        % identidad
\definecolor{acento}{HTML}{0F6A72}       % texto y linea de acento
\definecolor{acento2}{HTML}{0B5057}      % acento fuerte, enfasis
\definecolor{acentosuave}{HTML}{E3EFF0}  % fondo de acento
\definecolor{tinta}{HTML}{1E2829}        % texto principal
\definecolor{tinta2}{HTML}{49585A}       % texto secundario
\definecolor{tinta3}{HTML}{5C6B6D}       % texto terciario, notas
\definecolor{linea}{HTML}{DCE2E2}        % hairline
\definecolor{fondo}{HTML}{F4F6F6}        % fondo de pagina
\definecolor{superficie}{HTML}{EAEFF0}   % superficie elevada
\definecolor{bien}{HTML}{2F6B4C}         % estado favorable
\definecolor{aviso}{HTML}{8A5C17}        % estado en curso
\definecolor{mal}{HTML}{9C3A1F}          % estado que impide

% Clasificacion, misma escala que el reporte de lotes.
\definecolor{clasei}{HTML}{17563E}
\definecolor{claseii}{HTML}{2C5F9E}
\definecolor{claseiii}{HTML}{C08322}
\definecolor{claseiv}{HTML}{879596}

\setbeamercolor{normal text}{fg=tinta, bg=white}
\setbeamercolor{background canvas}{bg=white}
\setbeamercolor{frametitle}{fg=tinta, bg=white}
\setbeamercolor{framesubtitle}{fg=tinta2}
\setbeamercolor{structure}{fg=acento}
\setbeamercolor{alerted text}{fg=acento2}

\setbeamerfont{frametitle}{size=\large, series=\bfseries}
\setbeamerfont{framesubtitle}{size=\scriptsize, series=\mdseries}

% --- Rejilla fija ------------------------------------------------------------
% Todas las laminas comparten el mismo margen y la misma altura util, para que
% las proporciones no bailen entre una y otra.
\newlength{\margenlat}\setlength{\margenlat}{1.0cm}
\newlength{\altoutil}\setlength{\altoutil}{0.615\paperheight}

% Sin partir palabras: en laminas con columnas estrechas el corte silabico
% ensucia la lectura y en una presentacion no compensa.
\hyphenpenalty=10000
\exhyphenpenalty=10000
\setbeamersize{text margin left=\margenlat, text margin right=\margenlat}

% --- Grosores del sistema ----------------------------------------------------
\newlength{\anchofilete}\setlength{\anchofilete}{1.8cm}
\newlength{\altofilete}\setlength{\altofilete}{1.4pt}
\newlength{\hairline}\setlength{\hairline}{0.6pt}

% Divisor de marca. Es la unica linea gruesa de toda la presentacion.
\newcommand{\reglamarca}[1][acento]{\textcolor{#1}{\rule{\anchofilete}{\altofilete}}}

% Hairline de ancho libre, para separar bloques dentro de una lamina.
\newcommand{\filete}[1][\linewidth]{\textcolor{linea}{\rule{#1}{\hairline}}}

% Bandera moderada. \raggedright deja estiramiento infinito al margen derecho y
% en columna estrecha produce lineas de longitudes muy dispares. Acotando el
% estiramiento a dos cuadratines las lineas se igualan bastante sin justificar,
% que es lo que el sistema no quiere porque abre rios y partiria palabras.
\newcommand{\bandera}{\rightskip=0pt plus 2em\relax\parindent=0pt}

% Rotulo breve sobre un bloque. Siempre en versales y en gris terciario.
\newcommand{\rotulo}[2][tinta3]{{\scriptsize\textcolor{#1}{\MakeUppercase{#2}}}}

% --- Isotipo: siete barras de altura creciente, extremos redondeados ---------
\newcommand{\isotipo}[2][marca]{%
  \begin{tikzpicture}[y=#2, x=#2]
    \foreach \x/\y/\h in {0/0.270/0.460, 0.152/0.215/0.570, 0.304/0.160/0.680,
                          0.457/0.110/0.780, 0.609/0.060/0.880, 0.761/0.025/0.950,
                          0.909/0/1.0}
      {\fill[#1, rounded corners=0.0456*#2] (\x, \y) rectangle (\x+0.0913, \y+\h);}
  \end{tikzpicture}%
}

\newcommand{\marcaTexto}[1][marca]{%
  \raisebox{-0.28em}{\isotipo[#1]{0.62em}}\hspace{0.5em}%
  \textcolor{#1}{\textbf{\footnotesize MÉTODOS MIXTOS}}\,%
  \textcolor{tinta3}{\scriptsize CONSULTORES}%
}

% --- Titulo de lamina --------------------------------------------------------
% Titulo, subtitulo opcional y el filete de marca. Las medidas verticales son
% las mismas en todas las laminas: el cuerpo siempre arranca a la misma altura.
\makeatletter
\setbeamertemplate{frametitle}{%
  \vspace{0.45em}%
  \begin{beamercolorbox}[wd=\paperwidth, leftskip=\margenlat, rightskip=\margenlat]{frametitle}%
    % Sin \raggedright: su definicion pone \leftskip a cero y con el se llevaba
    % la sangria del propio beamercolorbox. Los titulos salian pegados al borde
    % del papel mientras el cuerpo arrancaba en el margen. Se compone en bandera
    % fijando las dos sangrias a mano.
    % El \rightskip se fija con \dimexpr: \margenlat es un registro de glue y
    % detras de un registro TeX no lee el "plus", que se imprimia como texto.
    \leftskip=\margenlat
    \rightskip=\dimexpr\margenlat\relax plus 1fil\relax
    \parindent=0pt%
    \usebeamerfont{frametitle}\insertframetitle%
    \ifx\insertframesubtitle\@empty\else%
      \\[0.25em]{\usebeamerfont{framesubtitle}\usebeamercolor[fg]{framesubtitle}\insertframesubtitle}%
    \fi%
  \end{beamercolorbox}%
  \vspace{0.5em}%
  % Sin \hspace: fuera del beamercolorbox ya rige el margen de lamina, de modo
  % que la sangria se sumaba dos veces y el filete arrancaba a 2 cm mientras el
  % titulo y el cuerpo arrancaban a 1 cm.
  \reglamarca%
  % Aire bajo el filete. Estaba en 0,85 em y dejaba una banda muerta entre el
  % filete y la figura justo donde la figura tenia que empezar, lo que obligaba a
  % encoger despues. Con 0,35 em la lamina recupera alto util y varias figuras
  % bajan de shrink.
  \vspace{0.35em}%
}
\makeatother

% --- Pie ---------------------------------------------------------------------
% Seccion a la izquierda, isotipo y numero a la derecha. Sin linea: el pie se
% separa por aire, no por trazo.
\setbeamertemplate{footline}{%
  \hbox{%
  \begin{beamercolorbox}[wd=0.5\paperwidth, ht=2.2ex, dp=1.6ex, leftskip=\margenlat]{}%
    \tiny\textcolor{tinta3}{\MakeUppercase{\insertsectionhead}}%
  \end{beamercolorbox}%
  \begin{beamercolorbox}[wd=0.5\paperwidth, ht=2.2ex, dp=1.6ex, rightskip=\margenlat]{}%
    \hfill\raisebox{-0.15em}{\isotipo[linea]{0.44em}}\hspace{0.5em}\tiny\textcolor{tinta3}{\insertframenumber}%
  \end{beamercolorbox}}%
}

% Sin justificar: en columnas estrechas la justificacion abre rios de espacio y
% parte palabras. La alineacion se marca dentro de cada minipage. No enganchar
% \frame con \apptocmd: rompe la lectura de [plain] y de los titulos.

\setbeamertemplate{itemize item}{\raisebox{0.14em}{\textcolor{acento}{\rule{0.4em}{0.4em}}}}
\setbeamertemplate{itemize subitem}{\raisebox{0.16em}{\textcolor{tinta3}{\rule{0.32em}{0.32em}}}}
\setlength{\leftmargini}{1.3em}
\setlength{\parskip}{0.4em}

% Tablas: hairline del sistema y aire suficiente entre filas.
\setlength{\heavyrulewidth}{\altofilete}
\setlength{\lightrulewidth}{\hairline}
\setlength{\aboverulesep}{0.5ex}
\setlength{\belowrulesep}{0.7ex}

% ============================================================================
% Piezas
% ============================================================================

% --- Separador de parte ------------------------------------------------------
% Lamina limpia: filete, rotulo de etapa, titulo grande y una linea de encuadre.
%
% Sin [plain]: con el, estas laminas salian sin pie y sin numero, y la numeracion
% visible saltaba de 5 a 7, de 8 a 10 y asi cinco veces. En un documento que se
% cita por lamina y sobre el que se van a mandar comentarios por escrito, eso
% confunde de inmediato. El pie queda, con su numero.
\newcommand{\parte}[3]{%
  \begin{frame}
    \vfill
    \hspace{\margenlat}\begin{minipage}{0.78\paperwidth}
      \raggedright
      \reglamarca\\[1.0em]
      \rotulo{#1}\\[0.35em]
      {\fontsize{26}{30}\selectfont\textcolor{tinta}{\textbf{#2}}}\\[0.7em]
      {\small\textcolor{tinta2}{#3}}
    \end{minipage}
    \vfill
  \end{frame}
}

% --- Lamina de afirmacion ----------------------------------------------------
% Una idea sostenida por un parrafo, con mucho aire alrededor.
\newcommand{\afirmacion}[2]{%
  \vspace{0.3em}
  \begin{minipage}{0.88\textwidth}
    \raggedright
    {\fontsize{15}{19}\selectfont\textcolor{tinta}{#1}}\\[0.9em]
    {\footnotesize\textcolor{tinta2}{#2}}
  \end{minipage}}

% --- Imagen a caja fija ------------------------------------------------------
% Tres cajas y solo tres, en vez de catorce alturas escritas a mano por lamina.
% La figura se dimensiona por la caja, de modo que dos laminas seguidas no
% cambian de peso. keepaspectratio hace que mande la restriccion mas estrecha.
%
%   \imagen        ancha, ocupa el ancho util entero
%   \imagenmedia   media, para la columna mayor de una lamina a dos columnas
%   \imagenalta    alta, para las capturas en retrato, donde manda la altura
\newcommand{\imagen}[1]{%
  \begin{center}\includegraphics[width=\textwidth, height=\altoutil, keepaspectratio]{#1}\end{center}}

\newcommand{\imagenmedia}[1]{%
  \includegraphics[width=\linewidth, height=0.52\paperheight, keepaspectratio]{#1}}

\newcommand{\imagenalta}[1]{%
  \includegraphics[width=\linewidth, height=0.66\paperheight, keepaspectratio]{#1}}

% --- Cifra -------------------------------------------------------------------
% Dato y su leyenda. La hairline superior las alinea cuando van varias en fila.
\newcommand{\cifra}[2]{%
  \begin{minipage}[t]{0.22\textwidth}
    \raggedright
    \filete\\[0.55em]
    {\color{acento}\fontsize{24}{26}\selectfont\textbf{#1}}\\[0.3em]
    {\scriptsize\color{tinta2} #2}
  \end{minipage}}

% Cifra de anclaje: la misma pieza en tamano de titular, para sostener una
% lamina de \afirmacion sin inventar contenido. Lleva la cifra que la frase ya
% menciona, no un dato nuevo.
\newcommand{\cifraancla}[3][0.30\textwidth]{%
  \begin{minipage}[t]{#1}
    \raggedright
    \filete\\[0.6em]
    {\color{acento}\fontsize{22}{25}\selectfont\textbf{#2}}\\[0.45em]
    {\scriptsize\color{tinta2} #3}
  \end{minipage}}

% Rotulo con hairline superior, sin cifra: para enumerar condiciones en fila.
\newcommand{\rotuloancla}[2][0.21\textwidth]{%
  \begin{minipage}[t]{#1}
    \raggedright
    \filete\\[0.6em]
    {\small\color{tinta}\textbf{#2}}
  \end{minipage}}

% --- Chip --------------------------------------------------------------------
% Etiqueta de estado. Nodo de tikz para que la caja se ajuste al texto y quede
% alineada con la linea base.
\newcommand{\chip}[2]{%
  \tikz[baseline=(c.base)]{%
    \node[rectangle, rounded corners=2pt, inner xsep=5pt, inner ysep=2.6pt,
          fill=#1!10, text=#1, font=\scriptsize\bfseries] (c) {#2};}}

% --- Cajas de diagrama -------------------------------------------------------
% El nodo ES la caja: el texto no puede desbordarla porque la caja se
% dimensiona a partir del texto. Dibujar rectangulo y texto por separado
% permite que el texto se salga, que es el defecto que esto corrige.
% El argumento del estilo es el ancho de texto. El inner sep de 7 pt es parte de
% la rejilla de los diagramas: cambiarlo estrecha el hueco de las flechas.
\tikzset{
  caja/.style={
    rectangle, rounded corners=2pt, draw=none,
    text width=#1, align=flush left, inner sep=7pt,
    anchor=north west, minimum height=3.55cm},
  cajarot/.style={
    rectangle, rounded corners=2pt, draw=linea, line width=\hairline,
    text width=#1, align=flush left, inner sep=7pt,
    anchor=north west, minimum height=3.55cm},
}

% Contenido tipico de una caja: rotulo, titulo y descripcion.
\newcommand{\cajatexto}[3]{%
  {\scriptsize\textcolor{acento}{\textbf{\MakeUppercase{#1}}}}\\[0.35em]
  {\small\textcolor{tinta}{\textbf{#2}}}\\[0.35em]
  {\scriptsize\textcolor{tinta2}{#3}}}

% --- Nota de fuente ----------------------------------------------------------
% Siempre al pie del contenido, en gris terciario y sin justificar.
% A 7 pt y no en scriptsize: la cita de fuente competia con el cuerpo de la
% lamina cuando iba debajo de una imagen. Es el minimo del sistema, el mismo del
% pie de lamina, de modo que sigue siendo legible proyectada.
\newcommand{\fuente}[1]{\vspace{0.3em}\par{\raggedright\fontsize{7}{8.4}\selectfont\color{tinta3} #1\par}}


% Los dos reportes se abren desde la propia lamina. El enlace es relativo: el PDF
% y los dos html viven en entregables/, de modo que funciona desde cualquier
% equipo mientras la carpeta viaje entera.
\hypersetup{colorlinks=true, urlcolor=acento, linkcolor=acento}

\makeatletter
\long\def\@textcolor#1#2#3{\protect\leavevmode{\color#1{#2}#3}}
\makeatother

\begin{document}

% ============================================================ 1. PORTADA
\begin{frame}[plain]
  \begin{tikzpicture}[remember picture, overlay]
    \node[anchor=east, inner sep=0pt]
      at ($(current page.east) - (0.55cm, 0)$)
      {\includegraphics[height=0.94\paperheight]{fig_mapa_portada.pdf}};
  \end{tikzpicture}%
  \vspace{0.2em}
  \begin{minipage}[c]{0.56\textwidth}
    \raggedright
    \isotipo{2.4em}\hspace{0.7em}%
    \raisebox{0.80em}{\begin{minipage}{6cm}
      \raggedright
      {\large\textcolor{marca}{\textbf{MÉTODOS MIXTOS}}}\\[0.1em]{\scriptsize\textcolor{tinta3}{C\,O\,N\,S\,U\,L\,T\,O\,R\,E\,S}}
    \end{minipage}}

    \vspace{1.1em}
    {\fontsize{26}{30}\selectfont\textcolor{tinta}{\textbf{Prospectos Solares}}}\\[0.45em]
    {\normalsize\textcolor{tinta2}{Encontramos los lotes donde cabe un proyecto solar\\
     y los entregamos medidos.}}

    \vspace{1.1em}
    \reglamarca

    \vspace{1.6em}
    {\small\textcolor{tinta}{Propuesta de servicio}}\\[0.55em]
    {\footnotesize\textcolor{tinta2}{Daniel Wiesner $\cdot$ Susana Yepes $\cdot$ Samuel Blanco}}\\[0.35em]
    {\scriptsize\textcolor{tinta3}{MÉTODOS MIXTOS CONSULTORES}}
  \end{minipage}
\end{frame}

% ============================================================ EL PROBLEMA
\section{El problema}

% --- 2. El problema, entero en una lamina ------------------------------------
% Eran dos: la de las cuatro condiciones y la de las tres entidades. Separadas,
% la primera enunciaba el problema y la segunda lo explicaba, y ninguna de las
% dos lo hacia evidente. Juntas dicen lo que hay que decir: el dato existe, y lo
% que no existe es el sitio donde esta reunido.
\begin{frame}{El dato existe. Reunido, no}
  \afirmacion{Cuatro condiciones tienen que coincidir sobre un mismo lote, y cada una la
  publica una entidad distinta.}
  {Ninguna responde por las demás, ninguna usa el mismo formato y ninguna comparte fecha
  de corte.}

  \vspace{1.9em}
  \hspace{0.01\textwidth}\begin{minipage}{0.96\textwidth}
    \raggedright
    \rotuloancla[0.22\linewidth]{Terreno apto y accesible}\hfill
    \rotuloancla[0.24\linewidth]{Capacidad libre en la red}\hfill
    \rotuloancla[0.24\linewidth]{Un uso del suelo que lo permita}\hfill
    \rotuloancla[0.22\linewidth]{Un dueño con quien negociar}
  \end{minipage}

  \vspace{1.9em}
  \hspace{0.01\textwidth}\begin{minipage}{0.96\textwidth}
    \bandera\small\textcolor{tinta2}{El catastro dice de quién es y cuánto mide, el
    ordenamiento del municipio si el uso lo permite, el operador de red si hay dónde
    conectarse.}
  \end{minipage}
\end{frame}

% ============================================================ LA PROPUESTA
\section{La propuesta}

% --- 3 ---
% El diagrama antes tenia los seis pasos de la version tecnica. Al cliente no le
% importa que artefacto viaja entre uno y otro: le importa que cada tramo entrega
% algo y que el ultimo entrega la ficha. Los tres primeros pasos son aqui una
% sola caja, «el modelo elige donde mirar».
\begin{frame}{Cómo trabajamos}
  \vfill
  \begin{center}
    \includegraphics[width=\textwidth, height=0.66\paperheight, keepaspectratio]{fig_cadena_comercial.pdf}
  \end{center}
  \vfill
\end{frame}

% --- 4. La unidad de trabajo, en grande --------------------------------------
\begin{frame}{La unidad de trabajo}{Una celda de 5 por 5 kilómetros, medida entera y a distancia}
  \vspace{0.2em}
  \begin{center}
    \includegraphics[width=\textwidth, height=0.62\paperheight, keepaspectratio]{fig_grilla_ancha.pdf}
  \end{center}
\end{frame}

% --- 5. Que se le mide a cada celda -----------------------------------------
% Estaba en cuatro columnas de 1, 3, 2 y 1 elementos: el terreno ocupaba el
% triple que el recurso sin que eso quiera decir nada, y los siete criterios no
% se leian como siete. Ahora es una rejilla de cuatro por dos, un criterio por
% tarjeta y el resultado en la octava casilla.
%
% «Cobertura del suelo» y no «Cobertura apta del suelo». El agregado sigue
% existiendo en el reporte de celdas, de modo que el nombre largo no seria falso
% aqui, pero el cliente no lo va a encontrar en el reporte de lotes: alli se
% retiro el 26 de agosto y la cobertura se reporta clase a clase. La nota al pie
% dice exactamente eso, para que la lamina no prometa una variable que el
% entregable ya no tiene con ese nombre.
\begin{frame}{Qué se le mide a cada celda}
  \vspace{0.1em}
  \centerline{\includegraphics[width=\textwidth, height=0.665\paperheight, keepaspectratio]{fig_criterios_celda.pdf}}
  \fuente{En la celda, la cobertura entra como un solo porcentaje de suelo apto. En la
  ficha del lote se reporta clase por clase, con la cifra de cada una y sin combinarlas.}
\end{frame}

% ============================================================ EL RESULTADO
\section{El resultado}

% --- 6 ---
\begin{frame}{Dónde mirar primero}{Las cien celdas mejor calificadas, sobre la red de transmisión}
  \vspace{-2.3em}
  \begin{minipage}[t]{0.705\textwidth}
    \vspace{0pt}
    \centering
    \includegraphics[width=\linewidth, height=0.755\paperheight, keepaspectratio]{fig_mapa.pdf}
  \end{minipage}\hfill
  \begin{minipage}[t]{0.265\textwidth}
    \vspace{1.8em}
    \bandera\footnotesize\textcolor{tinta2}{Las candidatas no se reparten por el país:
    siguen la red. Donde no hay línea cerca no hay nada que conectar, por bueno que sea
    el terreno.

    \vspace{0.9em}
    La nota compara cada celda con los terrenos donde ya hay planta. Dónde se pone el
    corte, un 70 por ejemplo, se decide con el cliente y cambia la lista.}
  \end{minipage}
\end{frame}

% --- 7 ---
% Figura propia de esta version, no la compartida con el guion tecnico. Aquella
% mete ademas la hoja entera a 2,15 cm de ancho, que a ese tamano no se lee y
% ademas se come una quinta parte del ancho. Fuera esa, las dos piezas graficas
% crecen de 5,8 y 3,2 cm a 8,4 y 4,6 cm.
\begin{frame}{Cada celda llega medida}{Una ficha por celda, con lo que decide si vale la pena bajar al catastro}
  \vspace{0.1em}
  \centerline{\includegraphics[width=\textwidth, height=0.60\paperheight, keepaspectratio]{fig_ficha_celda_comercial.pdf}}
  \fuente{Celda 0008656, Montería, vereda La Manta: prioritaria, índice 75 de 100. La hoja
  sigue con la aptitud de la celda y sus siete criterios, la potencia estimada, la
  cobertura, el recurso solar, el entorno, la conexión a la red y las restricciones.}
\end{frame}

% --- 8. De la celda al lote -------------------------------------------------
% Viene de la version tecnica. Es la lamina que ensena el salto de escala sin
% una sola cifra: a la izquierda lo que publica el catastro, a la derecha lo
% mismo ya clasificado.
\begin{frame}{Dentro de la celda, el catastro}{Se bajan todos los lotes y cada uno recibe sus propios valores}
  \begin{minipage}[t]{0.492\textwidth}
    \vspace{0pt}\centering
    \includegraphics[width=\linewidth, height=0.475\paperheight, keepaspectratio]{fig_catastro_celda.pdf}\\[0.4em]
    {\scriptsize\textcolor{tinta3}{Lo que baja del catastro}}
  \end{minipage}\hfill
  \begin{minipage}[t]{0.492\textwidth}
    \vspace{0pt}\centering
    \includegraphics[width=\linewidth, height=0.475\paperheight, keepaspectratio]{fig_catastro_clasificado.pdf}\\[0.4em]
    {\scriptsize\textcolor{tinta3}{Los mismos lotes, ya clasificados}}
  \end{minipage}
  \fuente{Celda 0010961, La Unión, Sucre. Azul idóneo, ámbar con gestión, rojo no viable.}
\end{frame}

% --- 9 ---
\begin{frame}{Usted pone las condiciones}{Se filtra como en una hoja de cálculo, y la lista se rehace a la vista}
  \vspace{0.3em}
  \begin{center}
    \includegraphics[width=\textwidth, height=0.505\paperheight, keepaspectratio]{fig_filtro_cliente.pdf}
  \end{center}
  \vspace{0.2em}
  \hspace{0.01\textwidth}\begin{minipage}{0.96\textwidth}
    \bandera\footnotesize\textcolor{tinta2}{Los diez de arriba son una muestra: el visor
    trae setenta y cinco criterios y ninguno trae el umbral puesto. Solo son fijos el mínimo
    de dos hectáreas y las exclusiones de entrada, y cada lote que cae por ahí va en la lista
    de descartes con su motivo.}
  \end{minipage}
\end{frame}

% --- 10 ---
\begin{frame}{Y clasificado por lo que le falta}{No es una lista de terrenos buenos: es una lista con el siguiente paso escrito}
  \vspace{0.3em}
  \begin{center}
    \includegraphics[width=\textwidth, height=0.535\paperheight, keepaspectratio]{fig_clases_lote.pdf}
  \end{center}
  \fuente{El Instituto Geográfico Agustín Codazzi no publica la zonificación rural de cuatro
  de cada diez lotes, ni de más de la mitad de los Idóneos: ahí decide el certificado de uso
  del suelo que expide el municipio.}
\end{frame}

% --- 11 y 12. Lo que trae la ficha del lote ---------------------------------
% La pregunta que el cliente hace siempre es que informacion recibe. La respuesta
% es el indice de la ficha, entero. Una captura de la hoja A4 no sirve, porque su
% cuerpo llega a la lamina por debajo de 3 pt.
%
% Van DOS laminas y no una. En una sola, a cuatro por cuatro, la tarjeta quedaba
% en 3,2 por 1,5 cm y el texto se salia de seis de ellas: no es un defecto de
% ajuste, es que diecisiete rotulos con su descripcion no caben en una hoja de
% 16 por 9. Partidas, la tarjeta pasa a 4,3 por 1,9 cm y el texto entra medido.
%
% Diecisiete y no dieciseis: comprobado el 2 de septiembre de 2026 sobre la ficha
% real, contando los h2 que compone el visor. La decimosexta es la del
% certificado de tradicion y solo sale cuando el folio ya esta leido, y eso lo
% dice el subtitulo de la primera de las dos.
\begin{frame}{Todo lo que trae la ficha de un lote}{Dieciséis secciones en una página imprimible, y una más cuando el certificado ya está leído}
  \vspace{0.1em}
  \centerline{\includegraphics[width=\textwidth, height=0.665\paperheight, keepaspectratio]{fig_ficha_completa_1.pdf}}
\end{frame}

\begin{frame}{Todo lo que trae la ficha de un lote}{De la norma urbanística a la ruta de compra}
  \vspace{0.1em}
  \centerline{\includegraphics[width=\textwidth, height=0.665\paperheight, keepaspectratio]{fig_ficha_completa_2.pdf}}
\end{frame}

% --- 13 ---
% El camino juridico es el diferenciador comercial: casi nadie lo entrega. Estaba
% con el quinto eslabon marcado como en desarrollo y ya no lo esta: al 2 de
% septiembre de 2026 funciona de punta a punta y esta probado sobre el folio
% 366-35594, de cinco paginas y doce anotaciones.
%
% Figura propia de esta version. La compartida con el guion tecnico sigue
% marcando el quinto eslabon como pendiente, y ese guion no se toca aqui.
%
% Cifras: el certificado a 23.000 pesos y la Resolucion 2026-001726 del 29 de
% enero de 2026 salen de outputs/COSTO_IA_CERTIFICADOS.md, que las verifico
% contra la resolucion; los 129 pesos y los 22.477 tokens son la medicion de esa
% misma nota sobre el folio real; los diez dias habiles, de docs/PREDIOS.md; los
% once riesgos, de la lista RIESGOS de predios/certificados.py.
\begin{frame}{Hasta el dueño y el folio}{Del código del catastro al certificado de tradición leído}
  \vspace{0.3em}
  \centerline{\includegraphics[width=\textwidth, height=0.53\paperheight, keepaspectratio]{fig_cadena_registro_comercial.pdf}}
  \fuente{La petición de matrículas no tiene costo y se responde en diez días hábiles. El
  certificado de tradición cuesta 23.000 pesos por lote, según la Resolución 2026-001726
  de la Superintendencia de Notariado y Registro, del 29 de enero de 2026, y por eso solo
  se compra sobre la lista ya filtrada. Leerlo cuesta unos 129 pesos, medidos sobre un
  folio real de cinco páginas y doce anotaciones, y un certificado ya leído no se vuelve a
  analizar. El modelo corre en el proyecto de Google Cloud de la firma.}
\end{frame}

% --- 14 ---
\begin{frame}{Qué recibe}{Tres entregables, y el procedimiento documentado para repetirlo}
  \hspace{0.01\textwidth}\begin{minipage}{0.96\textwidth}
    \bandera\footnotesize\textcolor{tinta2}{Todo llega como páginas web que se abren sin
    instalar nada y sin depender de nosotros.}
  \end{minipage}

  \vspace{0.6em}
  \begin{tikzpicture}[x=1cm, y=1cm]
    \node[cajarot=3.85cm, minimum height=2.85cm] (a) at (0, 0)
      {\cajatexto{Mapa del país}{Las celdas calificadas}{Las cien celdas candidatas, con
       su calificación y el motivo escrito. Se filtra y se descarga la selección.}};
    \node[cajarot=3.85cm, minimum height=2.85cm] (b) at (4.72, 0)
      {\cajatexto{Visor de lotes}{La lista corta}{Cada lote sobre imagen satelital, con
       sus medidas, su clase y los filtros en sus manos.}};
    \node[cajarot=3.85cm, minimum height=2.85cm] (c) at (9.44, 0)
      {\cajatexto{Ficha por lote}{El expediente}{Una página imprimible por lote, con lo
       que hay que gestionar y ante qué entidad.}};
  \end{tikzpicture}

  \vspace{0.7em}
  \hspace{0.01\textwidth}\begin{minipage}{0.96\textwidth}
    {\footnotesize\textcolor{tinta2}{Abrir ahora:
    \href{reporte_grillas.html}{\textbf{reporte de celdas}} $\cdot$
    \href{reporte_predios.html}{\textbf{visor de lotes}}}}
  \end{minipage}
\end{frame}

% ============================================================ CIERRE
\section{Cómo empezamos}

% --- 15 ---
% La columna derecha decia «Lo que entregamos primero: el mapa del pais con las
% zonas calificadas. Sobre el usted marca donde bajar al detalle. Nada sustituye
% la verificacion en campo». Vendia menos de lo que el producto hace: no hace
% falta que el cliente marque nada primero, y el mapa no es una entrega parcial a
% la espera de instrucciones. Con que diga sus condiciones de filtrado se le
% entrega el reporte entero, que es lo que se acaba de ensenar en once laminas.
%
% La salvaguarda de la verificacion en campo se queda, porque es cierta, pero
% baja al pie: puesta como tercer renglon de la columna parecia la conclusion de
% la lamina, y la conclusion es la entrega.
\begin{frame}{Cómo empezamos}
  \vspace{0.6em}
  \begin{minipage}[t]{0.47\textwidth}
    \raggedright
    \rotulo[acento]{Lo que necesitamos de usted}\\[0.7em]{\footnotesize\textcolor{tinta2}{La escala del proyecto que busca.\\[0.30em]
    Las zonas del país que le interesan.\\[0.30em]
    Sus condiciones de descarte, si ya las tiene.}}
  \end{minipage}\hfill
  \begin{minipage}[t]{0.47\textwidth}
    \raggedright
    \rotulo[acento2]{Lo que entregamos}\\[0.7em]{\footnotesize\textcolor{tinta2}{El mapa del país con las zonas calificadas,\\[0.30em]
    la lista corta de lotes y la ficha de cada uno.\\[0.30em]
    Suyos, para usarlos cuando quiera y sin nosotros.}}
  \end{minipage}

  \vspace{1.5em}
  \hspace{0.01\textwidth}\begin{minipage}{0.96\textwidth}
    \filete\\[0.9em]
    {\small\textcolor{tinta}{\textbf{Díganos qué busca y le entregamos el país ya consultado, lote por lote.}}}\\[0.45em]
    {\footnotesize\textcolor{tinta2}{Métodos Mixtos Consultores $\cdot$ Véalos funcionando:
    \href{reporte_grillas.html}{\textbf{reporte de celdas}} $\cdot$
    \href{reporte_predios.html}{\textbf{visor de lotes}}}}
  \end{minipage}

  \vspace{0.9em}
  \hspace{0.01\textwidth}\begin{minipage}{0.96\textwidth}
    \fuente{El trabajo se hace sobre cartografía y registros públicos: nada de lo que
    entregamos sustituye la verificación en campo.}
  \end{minipage}
\end{frame}

\end{document}
